\scanpage{10}
\mspage{III.6}
THEOREM III.6. If $\mathscr{L}$ is a locally separable $\handwritten{\mathscr{A}}$-algebra, and $\handwritten{\mathscr{R}}$ is a commutative $\handwritten{\mathscr{A}}$-algebra, then $\mathscr{L}\otimes_{\handwritten{\mathscr{A}}}\handwritten{\mathscr{R}}$ is a locally separable $\handwritten{\mathscr{R}}$-algebra.

\uline{Proof:} We know\handwritten{\unsure{,|$\swarrow$}} from [4, Corollary 1.6] that $(\mathscr{L}\otimes_{\handwritten{\mathscr{A}}}\handwritten{\mathscr{R}})_{\mathrm x}=\mathscr{L}_{\mathrm x}\otimes_{\handwritten{\mathscr{A}}_{\mathrm x}}\handwritten{\mathscr{R}}_{\mathrm x}$ is a central separable $\handwritten{\mathscr{R}}_{\mathrm x}$-algebra for every $\mathrm x$ in $\mathrm X$. Thus we need only show that $\mathscr{L}\otimes_{\handwritten{\mathscr{A}}}\handwritten{\mathscr{R}}$ has finite presentation as an $\handwritten{\mathscr{R}}$-algebra. For every $\mathrm x$ in $\mathrm X$ there is a neighborhood $\mathrm U$ of $\mathrm X$ and an exact sequence
\diagram{\begin{tikzcd}[ampersand replacement=\&,column sep=small]
0\handwritten{\arrow[r]}\& \handwritten{\mathscr{P}}\handwritten{\arrow[r]}\& \overset{\mathrm n}{\oplus}\,\handwritten{\mathscr{A}}|\mathrm U\handwritten{\arrow[r]}\& \mathscr{L}|\mathrm U\handwritten{\arrow[r]}\& 0
\end{tikzcd}}
where $\handwritten{\mathscr{P}}$ is an $\handwritten{\mathscr{A}}|\mathrm U$-module of finite type. Then the sequence $\handwritten{\mathscr{P}}\otimes_{\handwritten{\mathscr{A}}|\mathrm U}\handwritten{\mathscr{R}}|\mathrm U\handwritten{\longrightarrow}\overset{\mathrm n}{\oplus}\,\handwritten{\mathscr{A}}|\mathrm U\otimes_{\handwritten{\mathscr{A}}|\mathrm U}\handwritten{\mathscr{R}}|\mathrm U\handwritten{\longrightarrow} \mathscr{L}|\mathrm U\otimes_{\handwritten{\mathscr{A}}|\mathrm U}\handwritten{\mathscr{R}}|\mathrm U\handwritten{\longrightarrow} 0$ is exact, and $\handwritten{\mathscr{P}}\otimes_{\handwritten{\mathscr{A}}|\mathrm U}\handwritten{\mathscr{R}}|\mathrm U$ is an $\handwritten{\mathscr{R}}|\mathrm U$-module of finite type. Thus $\mathscr{L}\otimes_{\handwritten{\mathscr{A}}}\handwritten{\mathscr{R}}$ is finitely presented as an $\handwritten{\mathscr{R}}$-module, and therefore locally separable.

We recall now the notion of the commutant of a subring in a given ring, as described in [4]. This notion, generalized to sheaves, will be of use in the next section.

For any ring $\Lambda$ with center $\mathrm A$, and any two-sided $\Lambda$-module $\mathrm M$, $\mathrm{Hom}_{\Lambda^{\mathrm e}}(\Lambda,\mathrm M)$ is naturally isomorphic as an $\mathrm A$-module with $\mathrm M^\Lambda$, the $\mathrm A$-submodule of $\mathrm M$ consisting of all $\mathrm m$ in $\mathrm M$ such that $\mathrm{xm}=\mathrm{mx}$ for all $\mathrm x$ in $\Lambda$. If $\Gamma$ is a ring having center $\mathrm A$, and $\Lambda$ is an $\mathrm A$-subalgebra
% archive sheet number: 9
